% !PRINTED 628
\weberpage{628}
%<<CONT
% [sic] The source defines the resolvent as F(v)=0 but subsequently prints F(z) here; notation retained source-faithfully.
lichen Theiler, als die identische Permutation. Denn jeder gemeinsame Theiler aller dieser Gruppen müsste ein Normaltheiler der symmetrischen und folglich auch der alternirenden Gruppe sein; da er aber nicht die alternirende Gruppe selbst sein kann, so muss er sich nach \S.\,177 auf die Einheitsgruppe reduciren.

Die Resolvente $F(v) = 0$ ist also nach der in \S.\,156 eingeführten Bezeichnung eine \gesperrt{Totalresolvente} von $f(x) = 0$, und ihre Gruppe muss von gleichem Grade sein, wie die Gruppe von $f(x) = 0$, deren Grad $1 . 2 . 3 . 4 . 5 = 120$ ist. Da $F(z)$ auch irreducibel ist, so ist die Permutationsgruppe unter den Indices der sechs Grössen $v$ transitiv. Sie ist vom Grade 120, während der Grad der symmetrischen Gruppe der Permutationen von sechs Ziffern $6 . 120$ ist. Damit haben wir den Satz:

\gesperrt{Die symmetrische Permutationsgruppe von sechs Ziffern hat einen transitiven Divisor vom Index 6.}

Um diese merkwürdige Gruppe, die wir mit $C$ bezeichnen wollen, zu finden, haben wir nur den Einfluss zu untersuchen, den die sämmtlichen 120 Permutationen der $x$ auf die $v$ oder auf die Grössen (11) des vorigen Paragraphen ausüben.

Wir haben nun früher gesehen (\S.\,153, 2.), dass die Transpositionen $(0, 1)$, $(0, 2)$, $(0, 3)$, $(0, 4)$ durch wiederholte Zusammensetzung die ganze symmetrische Gruppe der Permutationen von fünf Ziffern erzeugen. Es genügt also, den Einfluss dieser vier Transpositionen auf die Indices der $v$ zu ermitteln. Durch Anwendung einer Transposition $(0, 1)$ geht jeder der Ausdrücke (11), \S.\,181 in einen der Ausdrücke (12) über und umgekehrt, und wenn man für die vier genannten Transpositionen diese Aenderung ermittelt, so erhält man die folgenden vier erzeugenden Permutationen $\pi_1, \pi_2, \pi_3, \pi_4$ der Gruppe $C$

\[
\begin{aligned}
(0, 1)&: \quad \pi_1 = (1, 3) (2, 5) (4, 6) \\
(0, 2)&: \quad \pi_2 = (1, 4) (2, 3) (5, 6) \\
(0, 3)&: \quad \pi_3 = (1, 5) (2, 6) (3, 4) \\
(0, 4)&: \quad \pi_4 = (1, 6) (2, 4) (3, 5),
\end{aligned}
\]

wo sich die in den $\pi$ vorkommenden Transpositionen $(1, 3) \ldots$ auf die Indices der $u$ oder der $v$ beziehen.

Unter den 120 Permutationen der Gruppe $C$ kommen unter anderen auch vor
%CONT>>
